\DeclareMathSymbol{:}{\mathbin}{operators}{"3A}


\usepackage{graphicx}
\usepackage{amstext}
\usepackage{wasysym}


\newcommand*{\Cdot}[1][2.5]{%
  \mathpalette{\CdotAux{#1}}\cdot%
}
\newdimen\CdotAxis
\newcommand*{\CdotAux}[3]{%
  {%
    \settoheight\CdotAxis{$#2\vcenter{}$}%
    \sbox0{%
      \raisebox\CdotAxis{%
        \scalebox{#1}{%
          \raisebox{-\CdotAxis}{%
            $\mathsurround=0pt #2#3$%
          }%
        }%
      }%
    }%
    % Remove depth that arises from scaling.
    \dp0=0pt %
    % Decrease scaled height.
    \sbox2{$#2\bullet$}%
    \ifdim\ht2<\ht0 %
      \ht0=\ht2 %
    \fi
    % Use the same width as the original \cdot.
    \sbox2{$\mathsurround=0pt #2#3$}%
    \hbox to \wd2{\hss\usebox{0}\hss}%
  }%
}


\newcommand{\Vars}{\mathit{Vars}}

\newcommand{\md}{\mathsf{md}}
\newcommand{\dmd}{\mathsf{dmd}}

% \renewcommand\merge[2]{%
%   \ooalign{\hfil$\vcenter{\hbox{$#1$}}$\hfil\cr
%     \hfil$\vcenter{\hbox{$#2$}}$\hfil}}

%\newcommand{\dyndiamond}{\raisebox{.4mm}{\merge\lozenge\Cdot}}
%\newcommand{\dynbox}{\raisebox{.4mm}{\merge\square\Cdot}}


\newcommand{\dyndiamond}[1]{{\fontsize{7}{7}\selectfont\tup{\raisebox{.6pt}{\textsf{#1}}}}}
\newcommand{\dynbox}[1]{{\fontsize{7}{7}\selectfont[\raisebox{.6pt}{\textsf{#1}}]}}

\newcommand\dswap{\dyndiamond{lsw}}
\newcommand\bswap{\dynbox{lsw}}
\newcommand\msw{\dyndiamond{lsw}}
\newcommand\mgsw{\dyndiamond{gsw}}

\newcommand\sab{\dyndiamond{gsb}}
\newcommand\sabbox{\dynbox{gsb}}
\newcommand\msb{\dyndiamond{lsb}}
\newcommand\mgsb{\dyndiamond{gsb}}

\newcommand{\lsab}{\dyndiamond{lsb}}
\newcommand{\lsabbox}{\dynbox{lsb}}

\newcommand{\bridge}{\dyndiamond{lbr}}
\newcommand{\bridgebox}{\dynbox{lbr}}
\newcommand\mbr{\dyndiamond{lbr}}
\newcommand\mgbr{\dyndiamond{gbr}}


\newcommand\dgswap{\dyndiamond{gsw}}
\newcommand\bgswap{\dynbox{gsw}}

\newcommand{\gbridge}{\dyndiamond{gbr}}
\newcommand{\gbridgebox}{\dynbox{gbr}}

\newcommand{\trbr}{\Circle^+}
\newcommand{\kn}{\textcircled{\small{k}}}
\newcommand{\re}{\textcircled{r}}

\newcommand{\qbfmodels}{\models_{\mathrm{qbf}}}

\newcommand{\iml}{\ensuremath{\mathsf{IML}}\xspace}
\newcommand{\bml}{\ensuremath{\mathcal{ML}}\xspace}
\newcommand{\bhl}{\ensuremath{\mathcal{HL}}\xspace}
\newcommand{\swl}{{\bml(\dswap)}}
\newcommand{\gsl}{{\bml(\sab)}}
\newcommand{\lsl}{{\bml(\lsab)}}
\newcommand{\brl}{{\bml(\bridge)}}
\newcommand{\gswl}{{\bml(\dgswap)}}
\newcommand{\gbrl}{{\bml(\gbridge)}}

% ---Memory Logics
\newcommand{\mlo}{{\cal ML}(\re,\kn)}

\newcommand{\trbrl}{{\bml(\trbr)}}

\newcommand{\fol}{{\mathcal{FOL}}}
%\newcommand{\sol}{{\mathcal{SOL}}}
\newcommand{\pow}{{\mathcal{P}}}

\newcommand{\down}{{\downarrow}}
\newcommand{\Hl}{\mathcal{H}}
\newcommand{\hlsd}{\Hl({:},\down)}
\newcommand{\hlsde}{\Hl({:},\down,\E)}

\newcommand{\A}{\textsf{A}}
\newcommand{\E}{\textsf{E}}
\newcommand{\PREF}{\mathsf{PREF}}
\newcommand{\DYN}{{\rm \sf DYN}\xspace}
\newcommand{\PROP}{{\rm \sf PROP}\xspace}
\newcommand{\VAR}{{\rm \sf VAR}\xspace}
\newcommand{\NOM}{{\rm \sf NOM}\xspace}
\newcommand{\FORM}{{\rm \sf FORM}\xspace}
\newcommand{\REL}{{\rm \sf REL}\xspace}
\newcommand{\ra}{\rightarrow}
\newcommand{\lra}{\leftrightarrow}
\newcommand{\Tr}{{\sf MTr}}
%\newcommand{\bisim}{\leftrightarroweq}
\renewcommand{\iff}{\mbox{iff}}
\newcommand{\eqrel}{\leftrightsquigarrow}
\renewcommand{\L}{{\cal L}}


%%%% Aca empiezan las macros viejas

\newcommand{\tup}[1]{\langle #1 \rangle}

\newcommand{\ttup}[1]{\langle\!\langle #1 \rangle\!\rangle}

\newcommand{\model}{\ensuremath{\mathcal{M}}\xspace}
\newcommand{\nodel}{\ensuremath{\mathcal{N}}\xspace}
\newcommand{\univ}{\ensuremath{\mathcal{U}}\xspace}
\newcommand{\uset}{U}
\newcommand{\img}{\textrm{im\ }}
\newcommand{\bisim}{\mathrel{\,\raisebox{.3ex}{$\underline{\makebox[.7em]{$\leftrightarrow$}}$}\,}}
\newcommand{\tbisim}{\bisim^T}

\newcommand{\cset}[1]{\{#1\}}
\newcommand{\csetsc}[2]{\{ #1 \mid #2\}}

\newcommand{\gM}{\mathcal{M}}
\newcommand{\gC}{\mathcal{C}}
\newcommand{\gL}{\mathcal{L}}


%%%%%%%%% de facu
% \theoremstyle{definition}
% \newtheorem{theorem}{Theorem}[section]
% \newtheorem{lemma}[theorem]{Lemma}
% \newtheorem{proposition}[theorem]{Proposition}
% \newcounter{coro}
% \newtheorem{corollary}[theorem]{Corollary}
% \newtheorem{conjecture}[theorem]{Conjecture}
% % \newtheorem{proof}{Proof}

% \newtheorem{obs}[theorem]{Observation}
% \newtheorem{definition}[theorem]{Definition}
% \newtheorem{construction}[theorem]{Construction}
% \newtheorem{nota}[theorem]{Notation}
% \newtheorem{example}[theorem]{Example}

%\newtheorem*{thm*}{Theorem} % sin numeracion
%\newtheorem*{defn*}{Definition}


\newcommand{\theq}[1]{\equiv_{#1}}

\newcommand{\elementary}{\ensuremath{EC_{\Delta}}\xspace}
\newcommand{\belementary}{\ensuremath{EC}\xspace}
\newcommand{\setcomp}[2]{\{{#1} : {#2}\}}

\newcommand{\diam}[1]{\langle #1 \rangle}
\newcommand{\ddiam}[1]{\ttup{#1}}
\newcommand{\ST}{\mathsf{STr}}

\newcommand{\HT}{\mathsf{HTr}}

%%%%%%%%%%%%%%%%%%%%%%%%%%%%%%%%%%%%%%
%---------Complexity Classes---------%
%%%%%%%%%%%%%%%%%%%%%%%%%%%%%%%%%%%%%%
\newcommand{\ntime}{\textsc{NTime}\xspace}
\newcommand{\ptime}{\textsc{P}\xspace}
\newcommand{\nptime}{\textsc{NP}\xspace}
\newcommand{\nlogspace}{\textsc{NLogSpace}\xspace}
\newcommand{\pspace}{\textsc{PSpace}\xspace}
\newcommand{\npspace}{\textsc{NPSpace}\xspace}
\newcommand{\exptime}{\textsc{ExpTime}\xspace}
\newcommand{\nexptime}{\textsc{NExpTime}\xspace}
\newcommand{\expspace}{\textsc{ExpSpace}\xspace}
\newcommand{\order}[1]{{\cal O} (#1)}
\newcommand{\C}{{\cal C}}

\newcommand{\wrv}{\tup{W,R,V}}
\newcommand{\wrvp}{\tup{W',R',V'}}

\newcommand{\EF}{\mathit{EF}}
\newcommand{\sbt}{\,\begin{picture}(-1,1)(-1,-3)\circle*{3}\end{picture}\ }


\newcommand{\sml}{\bml(\sdia)}
\newcommand{\shl}{\bhl(\sdia)}
